\documentclass[12pt]{article}
\usepackage[utf8]{inputenc}
\usepackage[T1]{fontenc}
\usepackage[margin=1in]{geometry}
\usepackage{setspace}
\usepackage{xcolor}
\usepackage[colorlinks=true, urlcolor=blue!60!black, linkcolor=blue!60!black]{hyperref}
\setlength{\parindent}{2em}
\setlength{\parskip}{1.25em}
\renewcommand{\baselinestretch}{1.0}
\title{\textit{Department of Homeland Security v. Thuraissigiam}, 591 U.S. 103 (2020)}
\author{Case Brief}
\date{}
\begin{document}
\maketitle

\section{Case Metadata}
\begin{itemize}
  \item \textbf{Case Name:} \textit{Department of Homeland Security v. Thuraissigiam}
  \item \textbf{Citation:} 591 U.S. 103 (2020) (No. 19-161); \href{https://www.supremecourt.gov/opinions/19pdf/19-161_g314.pdf}{SupremeCourt.gov Slip Opinion}; \href{https://www.law.cornell.edu/supremecourt/text/19-161}{Cornell LII Opinion}; \href{https://supreme.justia.com/cases/federal/us/591/19-161/}{Justia Opinion}
  \item \textbf{Dates:} Argued March 2, 2020; Decided June 25, 2020
  \item \textbf{Lower Courts:} U.S. District Court for the Southern District of California, 287 F. Supp. 3d 1077 (S.D. Cal. 2018) (dismissed for lack of jurisdiction); U.S. Court of Appeals for the Ninth Circuit, 917 F.3d 1097 (9th Cir. 2019) (reversed, holding 8 U.S.C. \S\ 1252(e)(2) unconstitutional)
  \item \textbf{Vote:} 7--2 (Reversed and remanded with instructions to dismiss)
  \item \textbf{Author:} Justice Alito (joined by Chief Justice Roberts and Justices Thomas, Gorsuch, and Kavanaugh)
  \item \textbf{Concurrences:} Justice Thomas (concurring); Justice Breyer (concurring in the judgment, joined by Justice Ginsburg)
  \item \textbf{Dissent:} Justice Sotomayor (dissenting, joined by Justice Kagan)
  \item \textbf{Constitutional Provisions:} Suspension Clause, U.S. Const. art. I, \S\ 9, cl. 2; Fifth Amendment Due Process Clause
\end{itemize}

\section{Facts and Procedural History}
Under the Illegal Immigration Reform and Immigrant Responsibility Act of 1996 (IIRIRA), certain noncitizens deemed ``applicants for admission'' who lack valid entry documents and have not been continuously physically present in the United States for two years are subject to expedited removal without a full removal hearing (\href{https://www.law.cornell.edu/supremecourt/text/19-161}{Cornell LII Opinion}). An applicant can exit expedited removal into standard removal proceedings by showing a ``credible fear of persecution''---a ``significant possibility'' of establishing asylum eligibility---before an asylum officer, with supervisory review and \textit{de novo} review by an immigration judge (8 U.S.C. \S\ 1225(b)(1)). Crucially, 8 U.S.C. \S\ 1252(e)(2) sharply restricts federal habeas corpus review of expedited removal orders to three narrow inquiries: (1) whether the petitioner is an alien, (2) whether the petitioner was ordered removed, and (3) whether the petitioner already holds lawful permanent resident, refugee, or asylee status. Federal courts may not review the merits of a negative credible-fear determination (\S\ 1252(a)(2)(A)(iii)).

In January 2017, Vijayakumar Thuraissigiam, a Sri Lankan national, crossed the U.S.--Mexico border unlawfully at night without inspection or entry documents. A Border Patrol agent apprehended him 25 yards north of the border, and DHS placed him in expedited removal (\href{https://www.law.cornell.edu/supremecourt/text/19-161}{Cornell LII Opinion}). During his credible-fear interview, Thuraissigiam stated that unidentified men had abducted and beaten him in Sri Lanka, but he did not link the attack to a protected ground. An asylum officer found no credible fear of persecution, a supervisor agreed, and an immigration judge affirmed \textit{de novo}.

Thuraissigiam filed a federal habeas petition in the Southern District of California, alleging for the first time persecution based on Tamil ethnicity and political opinion, claiming officials violated asylum procedures and applied the wrong legal standard, and seeking vacatur of his removal order and a new opportunity to apply for asylum. He did not request release from custody. The District Court dismissed for lack of jurisdiction under \S\ 1252(e)(2). The Ninth Circuit reversed, holding that \S\ 1252(e)(2) violated both the Suspension Clause and the Due Process Clause as applied, creating a split with the Third Circuit (\textit{Castro v. DHS}, 835 F.3d 422 (3d Cir. 2016)). The Supreme Court granted certiorari.

\section{Issue(s) Presented}
Whether 8 U.S.C. \S\ 1252(e)(2), which precludes federal habeas review of an expedited-removal negative credible-fear determination, violates the Suspension Clause (Art. I, \S\ 9, cl. 2) or the Fifth Amendment Due Process Clause as applied to a noncitizen apprehended shortly after crossing the border unlawfully.

\section{Holding}
No. In a 7--2 decision, the Supreme Court reversed and held that \S\ 1252(e)(2) violates neither the Suspension Clause nor the Due Process Clause as applied (\href{https://supreme.justia.com/cases/federal/us/591/19-161/}{Justia Opinion}). Because the writ of habeas corpus as understood in 1789 provided a means to secure release from unlawful executive detention---not to obtain administrative review or authorization to enter and remain in the country---the Suspension Clause did not guarantee judicial review of Thuraissigiam's asylum claim. Moreover, a noncitizen apprehended 25 yards inside the border stands ``on the threshold of initial entry'' and possesses only those procedural rights regarding admission that Congress has conferred by statute.

\section{Reasoning}
\subsection{Suspension Clause: The 1789 Core of Habeas Corpus}
Proceeding on the parties' shared premise that the Suspension Clause protects, at a minimum, ``the writ as it existed in 1789'' (\textit{INS v. St. Cyr}, 533 U.S. 289, 301 (2001)), Justice Alito emphasized that Blackstone, Justice Story, and Supreme Court precedent uniformly defined habeas corpus as a remedy to contest unlawful custody and secure simple release (\href{https://www.law.cornell.edu/supremecourt/text/19-161}{Cornell LII Opinion}). Thuraissigiam did not challenge the legality of his detention or seek simple release (which the Government was willing to effect by placing him on a flight to Sri Lanka); instead, he sought vacatur of his removal order and a new asylum hearing to gain lawful permission to remain in the United States---relief akin to an injunction or mandamus, and analogous to the claim rejected in \textit{Munaf v. Geren}, 553 U.S. 674 (2008).

Alito rejected three bodies of authority invoked by Thuraissigiam:
\begin{itemize}
  \item \textbf{Founding-Era Common-Law Cases:} Decisions such as \textit{Somerset v. Stewart} (1772) and early American deserting-sailor cases ordered simple release from unlawful physical confinement. Any resulting ability to stay in England or the United States was a collateral consequence of open-border immigration policy, not a habeas remedy granting admission.
  \item \textbf{``Finality Era'' Immigration Cases (1891--1952):} In \textit{Nishimura Ekiu v. United States}, 142 U.S. 651 (1892), and its progeny, federal courts reviewed legal questions in exclusion and deportation habeas cases not because the Suspension Clause compelled it---no majority opinion even cited the Clause---but because the Court read the 1891 Act's finality language narrowly to bar only factual review, leaving broad statutory habeas jurisdiction under the Habeas Corpus Act of 1867 intact.
  \item \textbf{\textit{Boumediene} and \textit{St. Cyr}:} \href{https://supreme.justia.com/cases/federal/us/553/723/}{\textit{Boumediene v. Bush}}, 553 U.S. 723 (2008), involved Guantanamo detainees seeking release from military custody, not admission to the United States, and \textit{St. Cyr} relied on statutory clear-statement canons for lawful residents facing deportation without approving habeas as a vehicle to seek entry.
\end{itemize}

\subsection{Due Process Clause: The ``Threshold of Initial Entry''}
Under more than a century of plenary-power precedent (\textit{Nishimura Ekiu}; \textit{United States ex rel. Knauff v. Shaughnessy}, 338 U.S. 537 (1950); \textit{Shaughnessy v. United States ex rel. Mezei}, 345 U.S. 206 (1953)), whatever procedure Congress authorizes is due process for a noncitizen seeking initial entry (\href{https://www.law.cornell.edu/supremecourt/text/19-161}{Cornell LII Opinion}). Alito rejected Thuraissigiam's claim that stepping 25 yards onto U.S. soil triggered full Fifth Amendment procedural protections: a noncitizen detained shortly after unlawful entry has not ``effected an entry'' (\textit{Zadvydas v. Davis}, 533 U.S. 678, 693 (2001)), and rewarding clandestine border crossings with greater constitutional rights than those afforded lawful port-of-entry applicants would create a perverse incentive. Because Thuraissigiam received the statutory credible-fear screening Congress provided, due process required nothing more.

\section{Concurring and Dissenting Opinions}
\begin{itemize}
  \item \textbf{Justice Thomas (Concurring):} Joined the majority and explored the original public meaning of the Suspension Clause, arguing that ``suspension'' historically meant a statute empowering the Executive to detain persons without bail or trial on mere suspicion of crime or dangerousness (\href{https://www.law.cornell.edu/supremecourt/text/19-161}{Cornell LII Opinion}). Because IIRIRA requires an inadmissibility finding and allows narrow habeas review, it is not a suspension.
  \item \textbf{Justice Breyer (Concurring in the Judgment, joined by Justice Ginsburg):} Criticized the majority for sweeping beyond the as-applied question and reaching due process (\href{https://www.law.cornell.edu/supremecourt/text/19-161}{Cornell LII Opinion}). Treating habeas as an ``adaptable remedy'' under \textit{Boumediene}, Breyer warned that the majority's rule could endanger long-term residents or citizens facing rogue executive abuses. He concurred narrowly because Thuraissigiam had never lived in the United States and his claims amounted to factual disputes and granular procedural quibbles that even finality-era habeas never reviewed.
  \item \textbf{Justice Sotomayor (Dissenting, joined by Justice Kagan):} Charged the majority with ``handcuff[ing]'' judicial oversight of executive detention and mischaracterizing Thuraissigiam's legal and procedural challenges as a bare plea for admission (\href{https://www.law.cornell.edu/supremecourt/text/19-161}{Cornell LII Opinion}). She argued that \textit{St. Cyr}, \textit{Boumediene}, \textit{Somerset}, and finality-era cases (as explained in \textit{Heikkila v. Barber}, 345 U.S. 229 (1953)) established that habeas tests the legality of executive removal orders and allows conditional release to cure errors. She also faulted the Court for extending the entry fiction to a person physically on U.S. soil without any clear limiting principle beyond 25 yards.
\end{itemize}

\section{Key Quotations}
\begin{quote}
``Habeas has traditionally been a means to secure release from unlawful detention, but respondent invokes the writ to achieve an entirely different end, namely, to obtain additional administrative review of his asylum claim and ultimately to obtain authorization to stay in this country.'' --- Justice Alito, \href{https://www.law.cornell.edu/supremecourt/text/19-161}{591 U.S. at 107}
\end{quote}

\begin{quote}
``The relief that a habeas court may order and the collateral consequences of that relief are two entirely different things.'' --- Justice Alito, \href{https://www.law.cornell.edu/supremecourt/text/19-161}{591 U.S. at 124}
\end{quote}

\begin{quote}
``For these reasons, an alien in respondent's position has only those rights regarding admission that Congress has provided by statute.'' --- Justice Alito, \href{https://www.law.cornell.edu/supremecourt/text/19-161}{591 U.S. at 140}
\end{quote}

\begin{quote}
``Today's decision handcuffs the Judiciary's ability to perform its constitutional duty to safeguard individual liberty and dismantles a critical component of the separation of powers. It will leave significant exercises of executive discretion unchecked in the very circumstance where the writ's protections `have been strongest.'\,'' --- Justice Sotomayor, dissenting, \href{https://www.law.cornell.edu/supremecourt/text/19-161}{591 U.S. at 161}
\end{quote}

\section{Doctrinal Significance and Modern Legacy (Through 2026)}
\textit{Thuraissigiam} marked a decisive turn toward strict 1789-originalism in Suspension Clause doctrine and bolstered Congress's plenary authority over immigration jurisdiction stripping:
\begin{itemize}
  \item \textbf{Expansion of 1789-Originalist Habeas Limits:} In \href{https://supreme.justia.com/cases/federal/us/599/21-857/}{\textit{Jones v. Hendrix}}, 599 U.S. 465 (2023), the Court extended \textit{Thuraissigiam}'s strict historical methodology beyond immigration, holding that Congress did not violate the Suspension Clause by barring federal prisoners from bringing statutory-interpretation innocence claims in successive 28 U.S.C. \S\ 2255 motions because 1789 habeas did not review final criminal judgments of courts of competent jurisdiction.
  \item \textbf{Preserving Core Habeas Review for Non-Border Removals (2025):} In \href{https://supreme.justia.com/cases/federal/us/604/24a931/}{\textit{Trump v. J.G.G.}}, 604 U.S. \_\_\_ (2025) (per curiam), and \href{https://supreme.justia.com/cases/federal/us/605/24a1007/}{\textit{A.A.R.P. v. Trump}}, 605 U.S. \_\_\_ (2025), the Court confirmed that noncitizens detained inside the United States and targeted for summary removal under the Alien Enemies Act are entitled under the Fifth Amendment to reasonable notice and an opportunity to challenge their confinement and removal via habeas corpus in the district of confinement---highlighting the boundary between summary interior removals and border-admission screening.
  \item \textbf{Boundary Disputes Over Interior Expedited Removal and Detention Through 2026:} As the Executive sought in 2025--2026 to expand expedited removal deep into the U.S. interior (up to the statutory two-year presence ceiling) and restrict border asylum access (see \href{https://supreme.justia.com/cases/federal/us/609/25-5/}{\textit{Mullin v. Al Otro Lado}}, 609 U.S. \_\_\_ (2026), addressing INA jurisdictional limits and ``arrives in the United States''), lower courts have policed \textit{Thuraissigiam}'s geographic and remedial limits. Notably, in \href{https://law.justia.com/cases/federal/appellate-courts/ca9/24-2801/24-2801-2026-09-21.html}{\textit{Padilla v. U.S. Immigration \& Customs Enforcement}}, No. 24-2801 (9th Cir. Sept. 21, 2026), the Ninth Circuit held that \textit{Thuraissigiam} did not extinguish the Fifth Amendment due process rights of noncitizens apprehended well within the United States---distinguishing them from entrants ``at the threshold of initial entry''---to seek bond hearings challenging prolonged detention.
\end{itemize}

\section{Pedagogical Critiques and Exam Hypotheticals}
\subsection{Key Critiques}
\begin{itemize}
  \item \textbf{Anachronistic 1789 Matching:} Critics echo Justice Sotomayor: demanding a 1789 common-law precedent for judicial review of asylum screenings asks for an impossibility when neither federal immigration restrictions nor asylum statutes existed in 1789, creating tension with \textit{Boumediene}'s functionalism.
  \item \textbf{Undefined ``Threshold of Entry'' Line:} By holding that 25 yards inside the border does not ``effect an entry'' without defining when territorial presence matures into constitutional personhood, the majority invited intense litigation over interior expedited removal up to two years after entry.
  \item \textbf{Unchecked Administrative Error:} Insulating credible-fear determinations from judicial review even for legal error risks \textit{refoulement} of bona fide refugees in high-volume, summary proceedings.
\end{itemize}

\subsection{Exam Hypotheticals and ``Opposite Side'' Contrast}
\textbf{Hypothetical 1 (Border-Adjacent Credible-Fear Denial):} A Guatemalan migrant crosses the Rio Grande without inspection and is apprehended by Border Patrol 100 yards north of the river two hours later. After an asylum officer and immigration judge deny her credible-fear claim, she files a \S\ 2241 habeas petition alleging the officer misapplied the statutory ``particular social group'' standard and seeking a new asylum interview. \textit{Analysis:} Under \textit{Thuraissigiam}, the district court must dismiss for lack of jurisdiction under 8 U.S.C. \S\ 1252(e)(2). Because she seeks a new opportunity to gain admission rather than simple release from unlawful custody, the Suspension Clause does not apply, and because she was apprehended immediately after unlawful entry at the threshold of the border, statutory procedure satisfies due process.

\textbf{Hypothetical 2 (``Opposite Side'' Contrast --- Core Habeas Release and Interior Due Process):} Suppose instead DHS arrests a noncitizen who has lived peacefully in Houston for eight years (or a U.S. citizen mistakenly flagged as an undocumented alien), detains him indefinitely without bond or notice of removal grounds, and invokes a new statute stripping all federal habeas review over executive detention orders. The detainee files a habeas petition demanding immediate release from unlawful physical confinement. \textit{Contrast:} Unlike \textit{Thuraissigiam}, this statute violates both the Suspension Clause and the Due Process Clause. The petitioner seeks \textit{simple release from unlawful executive detention}---the indisputable 1789 core of the Great Writ recognized in \textit{Thuraissigiam}, \textit{St. Cyr}, and \href{https://supreme.justia.com/cases/federal/us/553/723/}{\textit{Boumediene}}---and an established U.S. resident (or citizen) inside the country holds full Fifth Amendment due process protections against arbitrary detention and summary removal (\href{https://supreme.justia.com/cases/federal/us/604/24a931/}{\textit{Trump v. J.G.G.}} (2025); \href{https://law.justia.com/cases/federal/appellate-courts/ca9/24-2801/24-2801-2026-09-21.html}{\textit{Padilla v. ICE}} (9th Cir. 2026)).

\end{document}
